\scanpage{111}

\deleted{$G$ schéma \inserted{en groupes, \unsure{sur un anneau|mais encore|sur une base}} \unsure{prés.|qcq} \unsure{quelconque|quasi-sép.} et plat \unsure{de type fini} sur $k$ corps}

\deleted{Tout vectoriel $V$ sur $k$ où $G$ opère est $\varinjlim$ des sous-vectoriels $V'$ de dim finie sur $k$ stables par $G$ (SGAD $\mathrm{VI}_{\mathrm B}$ 11.2.) Donc la catégorie des}

$G$ schéma en groupes \inserted{affine et \deleted{plat} libre (pas néc. fini !)} sur $S$.

\deleted{On suppose l'une des conditions suivantes satisfaites}

\deleted{a) $G$ affine}

\deleted{b) $G$ qu. cpt \unsure{quasi sép|quasi-sép.} sur $S$ et plat sur $S$}

% slanted
\margin{N.B. il \unsure{suffit|suffirait} que $A(G)$ soit ``loc'' libre pour la topologie fpqc}

\deleted{Alors} Soit $\unsure{C|\mathcal C}_{\deleted{\unsure{f|\mathrm{qc}}}}(G)$ la catégorie des $S$-modules qu. coh. où $G$ opère, $C_f(G)$ la sous-catégorie formée des $\underline E$ qui sont de t.f. Alors

\[
C(G)\overset{\approx}{\longrightarrow}\operatorname{Ind} C_f(G)
\]

cf SGAD $\mathrm{VI}_{\mathrm B}$ 11.2.

Soit alors $f:T\longrightarrow S$ un morphisme \uline{affine}, défini par $\mathscr B=f_*(\underline{\mathcal O}_T)$, \deleted{et soit $f$ un} \inserted{Nous voudrions} \deleted{A \unsure{considérer|modules}} définir la catégorie des $G_T$-\deleted{E}Modules quasi-cohérents en termes de $C(G)$ et $\mathscr B$.

Or la donnée d'un $\underline{\mathcal O}_T$-module $\underline F$ qu. coh. revient à la donnée d'un $\underline E$ \inserted{$(=f_*(\underline F))$} qu. coh. sur $S$, sur lequel $\mathscr B$ opère. On constate de plus (sans hyp. sur $G$) que la donnée d'une opération
% archive sheet number: 110
